% \begin{figure}[t]
%     \centering
% \includegraphics[width=0.85\textwidth,trim = 1cm 0cm 1cm 0cm ]{img/compare.png}
% \caption{Comparison between ODE, SDE and Restart. \yx{TODO (1)}}   \label{fig:overview}
% \end{figure}

To sample from the aforementioned generative models, a prevalent approach employs general-purpose numerical solvers to simulate the corresponding differential equations. This includes Euler and Heun's 2nd method~\cite{Ascher1998ComputerMF} for ODEs (e.g., \Eqref{eq: diff-ode}), and Euler-Maruyama for SDEs (e.g., \Eqref{eq: diff-sde}). Sampling algorithms typically balance two critical metrics: (1) the quality and diversity of generated samples, often assessed via the Fréchet Inception Distance (FID) between generated distribution and data distribution~\cite{Heusel2017GANsTB}~(lower is better), and (2) the sampling time, measured by the number of function evaluations (NFE). Generally, as the NFE decreases, the FID score tends to deteriorate across all samplers. This is attributed to the increased discretization error caused by using a larger step size in numerical solvers.

However, as illustrated in \Figref{fig:main-line} and observed in previous works on diffusion models~\cite{Song2020ScoreBasedGM, Song2020DenoisingDI, Karras2022ElucidatingTD}, the typical pattern of the quality vs time curves behaves differently between the two groups of samplers, ODE and SDE. When employing standard numerical solvers, ODE samplers attain a decent quality with limited NFEs, whereas SDE samplers struggle in the same small NFE regime. However, the performance of ODE samplers quickly reaches a plateau and fails to improve with an increase in NFE, whereas SDE samplers can achieve noticeably better sample quality in the high NFE regime. This dilemma raises an intriguing question: \textit{Why do ODE samplers outperform SDE samplers in the small NFE regime, yet fall short in the large NFE regime?}

The first part of the question is relatively straightforward to address: given the same order of numerical solvers, simulation of ODE has significantly smaller discretization error compared to the SDE. For example, the first-order Euler method for ODE results in a local error of $O(\delta^2)$, whereas the first-order Euler-Maruyama method for SDEs yeilds a local error of $O(\delta^{\frac{3}{2}})$ (see \eg Theorem 1 of \cite{dalalyan2019user}), where $\delta$ denotes the step size. As $O(\delta^{\frac{3}{2}})\gg O(\delta^2)$, ODE simulations exhibit lower sampling errors than SDEs, likely causing the better sample quality with larger step sizes in the small NFE regime.

In the large NFE regime the step size $\delta$ shrinks and discretization errors become less significant for both ODEs and SDEs. In this regime it is the \textit{approximation error} — error arising from an inaccurate estimation of the ground-truth vector field by the neural network $s_\theta$ — starts to dominate the sampling error. We denote the discretized ODE and SDE using the learned field $s_\theta$ as $\bODE$ and $\bSDE$, respectively. In the following theorem, we evaluate the total errors from simulating $\bODE$ and $\bSDE$ within the time interval $[\tmin, \tmax] \subset [0,T]$. This is done via an upper bound on the Wasserstein-1 distance between the generated and data distributions at time $\tmin$. We characterize the accumulated initial sampling errors up until $\tmax$ by total variation distances. Below we show that the inherent stochasticity of SDEs aids in contracting these initial errors at the cost of larger additional sampling error in $[\tmin, \tmax]$. Consequently, SDE results in a smaller upper bound as the step size $\delta$ nears $0$ (pertaining to the high NFE regime).

\begin{theorem}[Informal]
\label{thm1_informal}
    Let $\tmax$ be the initial noise level and 
    % Let $\bx_{\tmax} = \by_{\tmax}$. Let
    % \begin{align*}
    %     \bx_{\tmin} = \bODE(\bx_{\tmax}, \tmax \to \tmin) \qquad \by_{\tmin} = \bSDE(\by_{\tmax}, \tmax \to \tmin)
    % \end{align*}
    $p_t$ denote the true distribution at noise level $t$. Let $p^{\bODE}_t, p^{\bSDE}_{t}$ denote the distributions of simulating $\bODE$, $\bSDE$ respectively. Assume that $\forall t \in  [\tmin, \tmax]$, $\lrn{x_t} < B/2$ for any $x_t$ in the support of $p_t$, $p^{\bODE}_t$ or $p^{\bSDE}_{t}$. Then
    {\footnotesize
    \begin{align*}
        & W_1(p^{\bODE}_{\tmin}, p_{\tmin}) \leq B \cdot TV\lrp{p^{\bODE}_{\tmax}, p_{\tmax}} + O(\delta+\eapprox) \cdot \lrp{\tmax - \tmin}\\
        & \underbrace{W_1(p^{\bSDE}_{\tmin}, p_{\tmin})}_{\textrm{total error}} \leq \underbrace{\lrp{1 - \lambda e^{-U}}  B \cdot TV(p^{\bSDE}_{\tmax}, p_{\tmax})}_{\textrm{upper bound on contracted error}} + \underbrace{O(\sqrt{\delta\tmax} + \eapprox) \lrp{\tmax - \tmin}}_{\textrm{upper bound on additional sampling error}}
    \end{align*}}
    In the above, {\footnotesize $U=BL_1/\tmin +L_1^2 \tmax^2 / \tmin^2$}, $\lambda < 1$ is a contraction factor, $L_1$ and $\eapprox$ are uniform bounds on $\lrn{ts_{\theta}(x_t,t)}$ and the approximation error $\lrn{t\nabla_x \log p_t(x) - ts_\theta(x,t)}$ for all $x_t,t$, respectively. $O()$ hides polynomial dependency on various Lipschitz constants and dimension.
\end{theorem}

We defer the formal version and proof of Theorem \ref{thm1_informal} to Appendix \ref{a:main_result}. As shown in the theorem, the upper bound on the total error can be decomposed into upper bounds on the \textit{contracted error} and \textit{additional sampling error}. {\footnotesize $TV(p^{\bODE}_{\tmax}, p_{\tmax})$} and {\footnotesize $TV(p^{\bSDE}_{\tmax}, p_{\tmax})$} correspond to the initial errors accumulated from both approximation and discretization errors during the simulation of the backward process, up until time $\tmax$. In the context of SDE, this accumulated error undergoes contraction by a factor of {\footnotesize $1 - \lambda e^{-BL_1/\tmin - L_1^2 \tmax^2 / \tmin^2}$} within $[\tmin, \tmax]$, due to the effect of adding noise. Essentially, the minor additive Gaussian noise in each step can drive the generated distribution and the true distribution towards each other, thereby neutralizing a portion of the initial accumulated error.

The other term related to additional sampling error includes the accumulation of discretization and approximation errors in $[\tmin, \tmax]$. Despite the fact that SDE incurs a higher discretization error than ODE ($O(\sqrt{\delta})$ versus $O(\delta)$), the contraction effect on the initial error is the dominant factor impacting the upper bound in the large NFE regime where $\delta$ is small. Consequently, the upper bound for SDE is significantly lower. This provides insight into why SDE outperforms ODE in the large NFE regime, where the influence of discretization errors diminishes and the contraction effect dominates. In light of the distinct advantages of SDE and ODE, it is natural to ask whether we can combine their strengths. Specifically, can we devise a sampling algorithm that maintains a comparable level of discretization error as ODE, while also benefiting from, or even amplifying, the contraction effects induced by the stochasticity of SDE? In the next section, we introduce a novel algorithm, termed \textit{Restart}, designed to achieve these two goals simultaneously.

% \xc{A few remarks: 
% \begin{enumerate}
%     \item Comparing SDE vs ODE, 
%     \begin{enumerate}
%         \item SDE has larger discretization error ($T\delta^{3/2}$ vs $T\delta^2$), but 
%         \item SDE has a contraction on the initial error.
%     \end{enumerate}
%     \item Comparing SDE VS Restart:
%     \begin{enumerate}
%         \item restart has $(1-\lambda)^K$.
%         \item restart has $\lambda$ compared to $\lambda e^{-BL/2}$ of SDE, 
%         \item restart has $\delta^2$ error compared to $\delta^{3/2}$
%     \end{enumerate}
% \end{enumerate}
% }



%\xc{we should have an experiment on a two-gaussian mixing showing that large noise leads to empirically better mixing, which will help address criticisms that we are comparing upper bounds.}
%\yx{sure. I guess we already have some old results regarding this. }